\section{Financial Markets}

\bd[Financial Asset]
A \textbf{financial asset} is a non-physical asset whose value is derived from a contract.
\ed

The contract that specifies the value of the financial asset is called a ``financial instrument''.

\bd[Financial Instrument]
A \textbf{financial instrument} is a contract that gives rise to a financial asset.
\ed

Sometimes the terms ``financial asset'' and ``financial instrument'' are used interchangeably, but it is important to
understand the difference between them. A financial asset is the result of a financial instrument, and a financial
instrument is a contract that gives rise to a financial asset. The financial instrument is the cause and the financial
asset is the effect. \v

Having defined financial instruments, we can now define a very special kind of market called a ``financial market''.

\bd[Financial Market]
A \textbf{financial market} is a market where people trade financial instruments at low transaction costs and at prices
that reflect supply and demand.
\ed

Financial markets can be divided into two categories: ``money markets'' and ``capital markets''. In these notes we will
not be discussing money markets at all, but we will be focusing on capital markets.

\subsection{Capital Markets}

\bd[Capital Market]
A \textbf{capital market} is a financial market in which long-term financial instruments are traded.
\ed

The most important financial instruments that are traded in capital markets are:
\bit
\item \textbf{Securities}: Financial instruments that represent ownership in a corporation or a creditor relationship
with a governmental body or a corporation.
\item \textbf{Derivatives}: Financial instruments that derive their value from the performance of an underlying asset,
index, or interest rate.
\eit

In this chapter and what follows we will be focusing on securities, and in the next chapter we will be focusing on
derivatives. \v

But before we start exploring securities, let's first have a look in the 4 main economic functions of capital markets
which are:
\bit
\item \textbf{Wealth Allocation} To facilitate the raising of capital by channeling wealth and allocating savings into
investments.
\item \textbf{Liquidity}: To allow the flexibility in timing consumption by providing liquidity of buying and selling
financial instruments.
\item \textbf{Risk Management}: To provide a mechanism for risk management including risk transfer, risk reduction
through diversification, and optimal risk sharing.
\item \textbf{Pricing}: To provide a mechanism for price discovery, i.e.\ the process of determining the price of an
asset in the market.
\eit

Hence, any time one engages in a capital market transaction they're altering the risk profile of their asset portfolios.
Subsequently, another way of thinking about capital markets is not necessarily the movement of capital, but the
alteration of those risks.

\section{Securities}

\bd[Security]
A \textbf{security} is a financial instrument that represents an ownership in a corporation or a creditor relationship
with a governmental body or a corporation.
\ed

\bd[Securities Market]
A \textbf{securities market} is a submarket of capital markets where people trade securities.
\ed

The main function of security markets is wealth allocation, i.e.\ to facilitate the transfer of funds from those who
have excess funds to those who need funds, by providing a mechanism for the buying and selling of securities.

\subsection{Market Participants}

Let's explore the 3 participants of security markets: investors, issuers, and financial intermediaries.

\bd[Investor]
An \textbf{investor} is a person or entity that commits capital with the expectation of receiving financial returns
generated by securities.
\ed

In traditional finance investors are considered to be just households, however in modern finance an investor can be any
institution who has capital to invest including pension funds investing group retirement plans, insurance companies
investing premiums, investment companies like mutual funds and unit trusts, sovereign wealth funds, endowments like
universities and non-profit oranizations, asset managers, and hedge funds. Its worths mentioning though that all these
institutional investors are ultimately investing the savings of households, hence in a way they act in behalf of the
households.

\bd[Issuer]
An \textbf{issuer} is a legal entity that develops, registers and sells securities for the purpose of financing its
operations.
\ed

In traditional finance issuers are considered to be just corporations, however in modern finance an issuer can be any
entity who needs to raise capital, including governments and governments agencies, municipalities, and supranational
entities. \v

Although issuers are the one who issue securities, most of the time they do not sell them directly to investors.
Instead,they hire financial intermediaries between them and the investors.

\bd[Financial Intermediary]
A \textbf{financial intermediary} is an institution that acts as the middleman between investors and issuers.
\ed

The main function of financial intermediaries is to facilitate the transfer of funds from those who have excess funds
to those who need funds, by providing a mechanism for buying and selling of securities.

\vspace{6pt}
\fig{cm}{0.57}
\vspace{6pt}

The most important financial intermediaries are comercial and investment banks.

\bd[Commercial Bank]
A \textbf{commercial bank} is a financial institution that provides deposit taking and commercial lending services.
\ed

\bd[Investment Bank]
An \textbf{investment bank} is a financial institution that provides securities trading, securities underwriting and
corporate advisory services.
\ed

Althoug there is a clear distinction between commercial and investment banks, in practice the distinction is not always
clear. In the past, commercial banks were not allowed to engage in investment banking activities, and investment banks
were not allowed to engage in commercial banking activities. However, in recent years, the distinction between the two
types of banks has become less clear, finding the two types of banks competing for similar services, and many banks
providing both types of services.

\subsection{Primary Market}

Capital markets are divided into two main types of ``submarkets'': the primary market and the secondary market.

\bd[Primary Market]
A \textbf{primary market} is the part of a security market where new securities are issued and sold for the first time.
\ed

New securities are issued and sold for the first time through initial public offerings (IPOs) or private placements.

\bd[Initial Public Offering (IPO)]
An \textbf{initial public offering} (\textbf{IPO}) is the first time that a company's securities are offered to the
public.
\ed

\bd[Private Placement]
A \textbf{private placement} is a sale of securities to pre-selected investors and institutions rather than on the
open market.
\ed

\subsection{Secondary Market}

\bd[Secondary Market]
A \textbf{secondary market} is the part of a security market where existing securities are bought and sold.
\ed

Secondary market is the trading of already outstanding securities. For the most part, they are non issuer
transactions. That is, the trade takes place between two different investors that are exchanging money for securities,
but it doesn't change the number of securities outstanding, nor does it provide the issuer with any additional
capital for investment. In simple words, there's no new money being raised and no wealth being channeled into
investment in the secondary market, however secondary markets are much larger than the primary markets, with the vast
majority of securities transactions (more than 95\%) taking place in the secondary market. \v

Although there are no issuers in secondary markets, there are financial intermediaries. The two primary roles of
financial intermediaries in the secondary market are to act either as ``market makers'' (also called ``dealers'') or
as ``brokers''.

\bd[Market Maker / Dealer]
A \textbf{market maker} or \textbf{dealer} is a financial intermediary that acts as principal by standing ready to
trade securities at all times at the publicly quoted price.
\ed

In simple words, a market maker is a financial intermediary in the business of trading securities for their own
account, and not on behalf of clients. Market makers take the opposite side of a trade from the customer, i.e.\ they
buy when a customer wants to sell and sell when a customer wants to buy. As it makes sense, this makes market makers
to have some capital in risk, due to the security inventory they have, and they need to be compensated for that. In
order to do so, market makers quote their own prices on securities both for buying (called ``bid price'') and for
selling (called ``ask price'' or ``offer price'').

\bd[Bid Price]
The \textbf{bid price} is the price at which a market maker is willing to buy a security.
\ed

\bd[Ask Price / Offer Price]
The \textbf{ask price}, or \textbf{offer price}, is the price at which a market maker is willing to sell a security.
\ed

It is important to note that both bid and ask prices do not reflect the true value of a security, but rather the market
maker's opinion of the security's value. The difference between the bid and ask price is called the ``spread'' and it
is the market maker's profit. The wider the spread, the more money the market maker makes.

\bd[Spread]
The \textbf{spread} is the difference between the ask price and the bid price:
\bse
\text{Spread} = \text{Ask Price} - \text{Bid Price}
\ese
\ed

The middle point between bid and ask prices is called the ``market price''.

\bd[Market Price]
The \textbf{market price} is the middle point between the bid and ask prices:
\bse
\text{Market Price} = \frac{\text{Bid Price} + \text{Ask Price}}{2}
\ese
\ed

The market price is the price at which a security is currently being traded in the secondary market. However, at the
market price, no one is able to neither buy nor sell a security. The bid and ask prices are the net prices at which a
market maker is willing to trade a security and that's the price that an investor will pay to buy a security or receive
to sell a security. \v

The difference between the market price and the bid and ask prices are called the ``mark-up'' and ``mark-down''.

\bd[Mark-Up]
A \textbf{mark-up} is the amount by which the ask price exceeds the market price:
\bse
\text{Mark-Up} = \text{Ask Price} - \text{Market Price}
\ese
\ed

\bd[Mark-Down]
A \textbf{mark-down} is the amount by which the bid price is less than the market price:
\bse
\text{Mark-Down} = \text{Market Price} - \text{Bid Price}
\ese
\ed

Saying the same thing in different words, market makers are eager to sell a security to an investor at market price
plus the mark-up, or buy a security from an investor at market price minus the mark-down. Mark-ups and mark-downs are
the profits market makers from selling and buying the securities. One can easily show that the spread is equal to the
sum of the mark-up and the mark-down:
\bse
\text{Spread} = \text{Mark-Up} + \text{Mark-Down}
\ese

Moving on, the second role of financial intermediaries in the secondary market is to act as brokers.

\bd[Broker]
A \textbf{broker} is a financial intermediary that brings buyers and sellers together but does not maintain an
inventory of securities.
\ed

A broker does not have capital at risk, and does not make money from the spread. Instead, brokers make money by
charging a commission for their services including compensation for time, effort and expenses of executing trade. The
commission is a fee that is charged for executing a trade. The commission can be a fixed amount, a percentage of the
trade, or a combination of both. \v

Secondary markets and their financial intermediaries provide several really important functions with the most important
one being liquidity. The fact that market makers make sure that securities can be bought and sold in the secondary
market provides investors with the ability to sell their securities and convert them into cash. This is important
because it means that investors can invest in securities without having to hold them until they mature. \v

Another important function of secondary markets is price discovery. The price of a security in the secondary market
provides information about the value of the security through the supply and demand for the security. This is important
because primary markets price things only once, at the time of the offering, and the price is set by the issuer,
without any input from the market and without having the possibility to change it, in a continuous changing world. \v

Up to this point we have discussed some of the generic characteristics of securities and the markets in which they are
traded, including the market participants and their roles, and the distinction between the primary and secondary market.
We will now switch gears and start discussing the two main types of securities.
\bit
\item \textbf{Debt / Fixed-Income Securities}: Securities that represent a creditor relationship with a governmental
body or a corporation.
\item \textbf{Equity Securities}: Securities that represent ownership in a corporation.
\eit

Let's start with debt/fixed-income securities.

\section{Debt / Fixed-Income Securities}

\bd[Debt / Fixed-Income Security]
A \textbf{debt security}, or \textbf{fixed-income security}, is a security that represents a creditor relationship with
a governmental body or a corporation.
\ed

In simple words, fixed-income securities are financial instruments that represent a loan made by an investor to the
issuer, making the investor a creditor of the issuer. The issuer of a fixed-income security is obligated to pay the
investor interest (hence the name ``fixed-income'') and to repay the initial loan amount at a later date, as it is
specified in the terms of repayment of the security.

\bd[Debt / Fixed-Income Markets]
A \textbf{debt market}, or \textbf{fixed-income market}, is a market where people trade debt securities.
\ed

Fixed-income markets are usually double or triple the size of equity markets, and are considered to be the most
important type of securities market. There are many different types of debt securities, but by far, the most important
ones are the so-called ``bonds''.

\subsection{Bonds}

\bd[Bond]
A \textbf{bond} is a fixed-income security that represents a loan made by an investor to a goverment, a govermental
agency or a corporate.
\ed

Bonds are the most important type of fixed-income securities. In fact, they are so important that the terms ``debt
security'', ``fixed-income security'', and ``bond'' are often used interchangeably, despite the fact that technically
they are not the same thing. Bonds are issued by governments, govermental agencies and corporations to raise capital.
Government and govermental agencies bonds (also called ``sovereign debts'') are usually the largest and most liquid
part of the fixed-income market, and are considered to be the safest type of fixed-income securities. As a result of
this, they have modest expected returns. \v

Corporate bonds are much more complex and difficult to categorize and analyze, however, they are usually riskier than
government and govermental agencies bonds, and as a result of this, they have higher expected returns. \v

When an investor buys a bond, they are lending money to the issuer in exchange for interest payments and the return of
the bond's principal when it matures.

\bd[Principal / Face Value / Par Value]
The \textbf{principal}, \textbf{face value}, or \textbf{par value} of a bond is the amount that the issuer agrees to
repay the bondholder at the maturity date.
\ed

\bd[Maturity Date]
The \textbf{maturity date} is the date on which the issuer of a bond agrees to repay to the bondholder the principal
of the bond.
\ed

The principal of the bond is usually paied out at its entirety at the end of the maturity date, however, there are
bonds that pay out the principal in installments, called ``amortizing bonds''.

\bd[Amortizing Bond]
An \textbf{amortizing bond} is a bond that pays out the principal in installments.
\ed

The interest payments on the principal that the issuer agrees to pay to the bondholder are called ``coupons''.

\bd[Coupon]
The \textbf{coupon} $C$ is the interest on the principal that the issuer agrees to pay to the bondholder.
\ed

\bd[Coupon Rate]
The \textbf{coupon rate} $r_C$ is the amount of interest due per period expressed as a percentage of the principal.
\ed

In simple words, the coupon is a regular payment the bond-holder receives, calculated by multiplying the coupon rate
by the principal of the bond and dividing by the number of coupon payments per year (compound) $n$:
\bse
C = \frac{r_C \cdot P}{n}
\ese

It is important to mention that coupon rates are usually fixed, but they can also be floating, i.e.\ they can change
over time.

\bd[Floating Coupon Rate]
A \textbf{floating coupon} is a coupon rate that can change over time.
\ed

Coupons are paid either all together at the end of the maturity date together with the principal (called ``zero-coupon
bonds''), or at regular time intervals, e.g.\ annually, semi-annually, quarterly, or monthly, depending on the
fixed-income market (called ``periodic coupon bonds'').

\bd[Zero-Coupon Bond]
A \textbf{zero-coupon bond} is a bond in which coupons are paied at their entirety at the end of the maturity date.
\ed

\bd[Periodic Coupon Bond]
A \textbf{periodic coupon bond} is a bond in which coupons are paied at regular time intervals.
\ed






As it makes sense, bonds carry a certain amount of risks that makes the future return of the bondholder uncertain:
\bit
\item \textbf{Default / Credit Risk}: The risk that the issuer will not be able to make the interest payments or repay
the principal at the maturity date.
\item \textbf{Interest Rate Risk}: The risk that the value of the bond will decrease due to changes in interest rates.
This risk can be spllited into two subcategories: ``price risk'' and ``reinvestment risk''.
\bit
\item \textbf{Market / Price Risk}: The risk that the value of the bond will decrease due to an increase in interest
rates.
\item \textbf{Reinvestment Risk}: The risk that the bondholder will not be able to reinvest the interest payments at
the same interest rate.
\eit
\item \textbf{Inflation Risk}: The risk that the value of the bond will decrease due to an increase in inflation.
\item \textbf{Liquidity Risk}: The risk that the bondholder will not be able to sell the bond at a fair price.
\eit

Of course, since bonds carry risks they also carry sources of returns:
\bit
\item \textbf{Coupons}: The interest payments that the bondholder receives from the issuer.
\item \textbf{Principal Capital Gains}: The increase in the face value of the bond due to a decrease in interest rates.
\item \textbf{Compound Interest}: The interest that the bondholder receives from reinvesting the interest payments.
\eit




\bd[Bond Price]
The \textbf{price} of a bond is the present value of the bond given the current yield that the bond is traded.
\ed

\bd[Bond Valuation]
The \textbf{valuation} of a bond is the process of determining what a bond is worth through analytic techniques.
\ed

To price a periodic coupon bond, we interpret it as a series of cash flows and discount them back to the present value.

\fig{cm3}{0.5}

The cash flows of a periodic coupon bond are the periodic coupons $C$ and the principal $P$ received at maturity. The
discount rate is the yield at maturity $r$ of the bond. By assuming $n$ compounding periods, the price of a bond is
given by the formula for the present value of a series of cash flows (\ref{eq:pv}):
\bse
PV = \sum_{i=1}^T \Big(\frac{C}{(1 + \frac{r}{n})^{n \cdot i}}\Big) + \frac{P}{(1 + \frac{r}{n})^{n \cdot T}}
\ese

where:
\bse
C = \frac{r_C \cdot P}{n}
\ese

Notice that in this case we have assumed a non-flotating coupon rate.



%\section{Equity Securities}
%
%\bd[Equity Security]
%An \textbf{equity security} is a security that represents ownership in a corporation.
%\ed
%
%Equity securities are also called ``stocks'' or ``shares''. When you buy a stock you are buying a piece of the
%company. The company is divided into a number of shares and each share represents a piece of the company. When you
%buy a share you become a shareholder of the company. The company is legally required to provide you with certain
%rights and privileges as a shareholder. These rights and privileges include the right to vote in the company's annual
%general meeting, the right to receive dividends, the right to receive a portion of the company's assets in case the
%company is liquidated, and others.